\subsection{Threats to validity}
\label{sec:threats}

\textbf{Construct validity.}  Our defects are injected by operators, not
collected from real rejected pull requests, so their distribution follows the
operators and the constructs that happen to occur near the patched code.  Some
operators leave local cues that a pattern matcher can exploit (for example
\texttt{except Exception: pass}); the static-analysis baseline measures how far
such cues go.  The plausibility and the defectiveness of the variants were not
validated by human annotators.

\textbf{Internal validity.}  The admission filter runs the PR's relevant test
files rather than the whole suite; a variant could in principle break a test
elsewhere in the repository, which would make the task easier for an execution
oracle but does not change what the judge sees.  Docker images were pulled
through a public registry mirror and are the images published with the tasks.

\textbf{Judges and prompting.}  All judges are small (1.5--4B), quantised
(4--8 bit) open models.  In the main condition they must answer without first
writing an explanation; the base paper lets judges reason before the verdict,
so our direct-verdict numbers are a lower bound for these models, and the
thinking condition measures how much reasoning recovers.

\textbf{Statistical conclusion validity.}  The benchmark is small, rare classes
have few pairs, and pairs built from the same pull request are correlated; we
therefore resample tasks rather than pairs in the bootstrap and report
intervals rather than significance stars.  We do not correct for multiple
comparisons; per-class and per-condition differences should be read as
exploratory.

\textbf{Selection.}  Pairs whose diff-only prompt exceeds the judges' token
budget, and tasks with solution leakage, are excluded from judging, which caps
the diff sizes the scaling analysis can observe.
